\subsection{From equivariant local features to sparse operator prediction}
\label{subsec:matrix_observable_background}

\texttt{Mandala} represents each operator $X \in \{H,S,D\}$ as a sparse family of
atom-pair blocks $X_{ij}^{\bm{L}}$, where $i$ and $j$ index atoms and
$\bm{L}$ labels periodic images. The block representation is converted into an
irrep representation by a fixed change of basis,
\begin{equation}
\bm{x}_{ij}^{\bm{L}} = Q_{a_i a_j} \, \mathrm{vec}\!\left(X_{ij}^{\bm{L}}\right),
\end{equation}
where $a_i$ and $a_j$ denote the atomic species of atoms $i$ and $j$, and
$Q_{a_i a_j}$ is the species-pair-specific block-to-irrep transform. The
inverse map reconstructs matrix blocks from predicted irrep vectors,
\begin{equation}
\mathrm{vec}\!\left(X_{ij}^{\bm{L}}\right)
=
Q_{a_i a_j}^{-1} \bm{x}_{ij}^{\bm{L}},
\end{equation}
which, in the orthonormal convention used internally, is implemented as a
transpose. In the code, these maps are materialized by the
\texttt{BlockIrrepMapper}, which is reused by the data pipeline, the model
heads, and the evaluation routines.

Once the network predicts a consistent set of sparse operators, observable
quantities are obtained directly from those operators rather than through a
separate surrogate model. For sparse block matrices, the trace of a product is
evaluated as
\begin{equation}
\Tr(AB)
=
\sum_{(\bm{L},i,j)}
\mathrm{tr}\!\left(
A_{ij}^{\bm{L}} B_{ji}^{-\bm{L}}
\right),
\end{equation}
which is exactly the sparse contraction implemented by \texttt{Mandala}. This gives
direct access to electron counts, energy observables, and
derivative observables such as forces and stress, all from the same predicted
operator set. The resulting workflow differs from direct scalar regression:
internal consistency between matrices and observables is built into the model
formulation rather than imposed only at the level of post hoc analysis.
